% BEGIN STRUCTURAL ARTIFACTS: mathematical-spaces
\begin{remark*}[Structural Cards]
The following cards use the structural presentation layer: each card records
the inherited arena, the new data or laws, and the role of the resulting
space.  These are blueprint artifacts, not formal definitions; formal
definitions remain authoritative when they are introduced.

\begin{structureblueprint}{Structural Card: Set}
\StructureMetadata{id=struct:space-set; kind=structure; structure=set}
\begin{structuresection}{TAXONOMY}
Primitive arena.
\end{structuresection}
\begin{structurenew}
A nonempty domain \(X\) of objects.
\end{structurenew}
\begin{structuretheory}
No nearness, size, operation, or distinguished family of subsets is specified.
\end{structuretheory}
\end{structureblueprint}

\begin{structureblueprint}{Structural Card: Family of Subsets}
\StructureMetadata{id=struct:space-family-of-subsets; kind=structure; structure=family-of-subsets}
\begin{structureinherits}
Set \(X\).
\end{structureinherits}
\begin{structurenew}
A chosen collection \(\mathcal A\subseteq\mathcal P(X)\).
\end{structurenew}
\begin{structuresection}{TAXONOMY}
Subset-selection structure.  No closure law is imposed at this level.
\end{structuresection}
\end{structureblueprint}

\begin{structureblueprint}{Structural Card: Topological Space}
\StructureMetadata{id=struct:space-topological; kind=structure; structure=topological-space}
\begin{structureinherits}
Set \(X\) with a distinguished family of subsets.
\end{structureinherits}
\begin{structurenew}
A topology \(\mathcal T\subseteq\mathcal P(X)\), read as the open subsets of
\(X\).
\end{structurenew}
\begin{structuretheory}
\(\varnothing,X\in\mathcal T\); arbitrary unions of members of
\(\mathcal T\) lie in \(\mathcal T\); finite intersections of members of
\(\mathcal T\) lie in \(\mathcal T\).
\end{structuretheory}
\begin{structuresection}{TAXONOMY}
Open-set structure for continuity, convergence, interior, closure, and
neighborhood arguments.
\end{structuresection}
\end{structureblueprint}

\begin{structureblueprint}{Structural Card: Metric Space}
\StructureMetadata{id=struct:space-metric; kind=structure; structure=metric-space}
\begin{structureinherits}
Set \(X\).
\end{structureinherits}
\begin{structurenew}
A distance function
\[
  d:X\times X\to[0,\infty).
\]
\end{structurenew}
\begin{structuretheory}
Identity of indiscernibles, symmetry, and the triangle inequality.
\end{structuretheory}
\begin{structureusesa}
Open balls \(B(x,r)=\{y\in X:d(x,y)<r\}\).
\end{structureusesa}
\begin{structuresection}{TAXONOMY}
Distance structure.  The metric induces a topology by declaring sets open
when they contain a metric ball around each of their points.
\end{structuresection}
\end{structureblueprint}

\begin{structureblueprint}{Structural Card: Vector Space}
\StructureMetadata{id=struct:space-vector; kind=structure; structure=vector-space}
\begin{structureinherits}
Set \(V\) and a scalar field \(\mathbb F\).
\end{structureinherits}
\begin{structurenew}
Vector addition \(+:V\times V\to V\), scalar multiplication
\(\cdot:\mathbb F\times V\to V\), and a zero vector \(0\in V\).
\end{structurenew}
\begin{structuretheory}
\((V,+,0)\) is an abelian group; scalar multiplication is unital and
associative; scalar multiplication distributes over vector addition and field
addition.
\end{structuretheory}
\begin{structuresection}{TAXONOMY}
Linear structure.  The space supports linear combinations before any length,
angle, or limit structure is chosen.
\end{structuresection}
\end{structureblueprint}

\begin{structureblueprint}{Structural Card: Normed Vector Space}
\StructureMetadata{id=struct:space-normed-vector; kind=structure; structure=normed-vector-space}
\begin{structureinherits}
Vector space \(V\) over \(\mathbb F\).
\end{structureinherits}
\begin{structurenew}
A norm \(\lVert\cdot\rVert:V\to[0,\infty)\).
\end{structurenew}
\begin{structuretheory}
\(\lVert v\rVert=0\) if and only if \(v=0\);
\(\lVert\alpha v\rVert=|\alpha|\,\lVert v\rVert\);
\(\lVert u+v\rVert\leq\lVert u\rVert+\lVert v\rVert\).
\end{structuretheory}
\begin{structureusesa}
The induced metric \(d(u,v)=\lVert u-v\rVert\).
\end{structureusesa}
\begin{structuresection}{TAXONOMY}
Linear metric structure.  Algebraic operations and metric notions become
compatible.
\end{structuresection}
\end{structureblueprint}

\begin{structureblueprint}{Structural Card: Inner Product Space}
\StructureMetadata{id=struct:space-inner-product; kind=structure; structure=inner-product-space}
\begin{structureinherits}
Real or complex vector space \(V\).
\end{structureinherits}
\begin{structurenew}
An inner product \(\langle\cdot,\cdot\rangle:V\times V\to\mathbb F\).
\end{structurenew}
\begin{structuretheory}
Positive definiteness, conjugate symmetry, and linearity in the chosen
argument.
\end{structuretheory}
\begin{structureusesa}
The induced norm \(\lVert v\rVert=\sqrt{\langle v,v\rangle}\).
\end{structureusesa}
\begin{structuresection}{TAXONOMY}
Linear geometry.  Length, angle, and orthogonality are available, and the
metric structure is derived from the inner product.
\end{structuresection}
\end{structureblueprint}

\begin{structureblueprint}{Structural Card: Banach Space}
\StructureMetadata{id=struct:space-banach; kind=structure; structure=banach-space}
\begin{structureinherits}
Normed vector space \((V,\lVert\cdot\rVert)\).
\end{structureinherits}
\begin{structurenew}
Completeness with respect to the norm-induced metric.
\end{structurenew}
\begin{structuretheory}
Every Cauchy sequence in \(V\) converges to a point of \(V\) in the norm.
\end{structuretheory}
\begin{structuresection}{TAXONOMY}
Complete normed linear space.  Approximation by Cauchy sequences remains
inside the space.
\end{structuresection}
\end{structureblueprint}

\begin{structureblueprint}{Structural Card: Hilbert Space}
\StructureMetadata{id=struct:space-hilbert; kind=structure; structure=hilbert-space}
\begin{structureinherits}
Inner product space \(H\).
\end{structureinherits}
\begin{structurenew}
Completeness with respect to the norm induced by the inner product.
\end{structurenew}
\begin{structuretheory}
Every Cauchy sequence in the induced norm converges to a point of \(H\).
\end{structuretheory}
\begin{structuresection}{TAXONOMY}
Complete inner product space.  Orthogonality and projection arguments take
place in a complete metric environment.
\end{structuresection}
\end{structureblueprint}

\begin{structureblueprint}{Structural Card: Measurable Space}
\StructureMetadata{id=struct:space-measurable; kind=structure; structure=measurable-space}
\begin{structureinherits}
Set \(X\) with a distinguished family of subsets.
\end{structureinherits}
\begin{structurenew}
A \(\sigma\)-algebra \(\mathcal M\subseteq\mathcal P(X)\), read as the
measurable subsets of \(X\).
\end{structurenew}
\begin{structuretheory}
\(X\in\mathcal M\); complements of measurable sets are measurable; countable
unions of measurable sets are measurable.
\end{structuretheory}
\begin{structuresection}{TAXONOMY}
Countably closed subset-selection structure.  It specifies which subsets are
legal inputs for size assignments.
\end{structuresection}
\end{structureblueprint}

\begin{structureblueprint}{Structural Card: Measure Space}
\StructureMetadata{id=struct:space-measure; kind=structure; structure=measure-space}
\begin{structureinherits}
Measurable space \((X,\mathcal M)\).
\end{structureinherits}
\begin{structurenew}
A measure \(\mu:\mathcal M\to[0,\infty]\).
\end{structurenew}
\begin{structuretheory}
\(\mu(\varnothing)=0\), and \(\mu\) is countably additive on pairwise disjoint
measurable sets:
\[
  \mu\!\left(\bigcup_{n=1}^{\infty}A_n\right)
  =
  \sum_{n=1}^{\infty}\mu(A_n).
\]
\end{structuretheory}
\begin{structuresection}{TAXONOMY}
Quantitative measurable structure.  The \(\sigma\)-algebra determines the
admissible subsets, and the measure assigns their size.
\end{structuresection}
\end{structureblueprint}
\end{remark*}
% END STRUCTURAL ARTIFACTS: mathematical-spaces
